\subsection{ELEMENTARY AUTOMORPHISMS}

Let \(\mathfrak{g}\) be a Lie algebra. Denote its group of automorphisms by \(\operatorname{Aut}(\mathfrak{g})\) . If \(x \in \mathfrak{g}\) and if \(\operatorname{ad} x\) is nilpotent, \(e^{\operatorname{ad} x} \in \operatorname{Aut}(\mathfrak{g})\) (Chap. I, §6, no. 8).

DEFINITION 1. A finite product of automorphisms of g of the form \(e^{ad x}\) with ad x nilpotent is called an elementary automorphism of g. The group of elementary automorphisms of g is denoted by \(\operatorname{Aut}_{e}(\mathfrak{g})\) .

If \(u \in \operatorname{Aut}(\mathfrak{g})\) , \(ue^{\operatorname{ad}x}u^{-1} = e^{\operatorname{ad}u(x)}\) . It follows that \(\operatorname{Aut}_{e}(\mathfrak{g})\) is a normal subgroup of \(\operatorname{Aut}(\mathfrak{g})\) . If k = R or C, \(\operatorname{Aut}_{e}(\mathfrak{g})\) is contained in the group \(\operatorname{Int}(\mathfrak{g})\) of inner automorphisms of g (Chap. III, §6, no. 2, Def. 2).

* In the general case, \(\mathrm{Aut}_e(\mathfrak{g})\) is contained in the identity component of the algebraic group \(\mathrm{Aut}(\mathfrak{g})\) .*

Lemma 1. Let \(\mathrm{V}\) be a finite dimensional vector space, \(\mathfrak{n}\) a Lie subalgebra of \(\mathfrak{a} = \mathfrak{gl}(\mathrm{V})\) consisting of nilpotent elements.

\begin{enumerate}
\def\labelenumi{(\roman{enumi})}
\item
  The map \(x \mapsto \exp x\) is a bijection from \(\mathfrak{n}\) to a subgroup \(N\) of \(\mathbf{GL}(V)\) consisting of unipotent elements (Chap. II, §6, no. 1, Remark 4). We have \(\mathfrak{n} = \log(\exp \mathfrak{n})\) . The map \(f \mapsto f \circ \log\) is an isomorphism from the algebra of polynomial functions on \(\mathfrak{n}\) to the algebra of restrictions to \(\mathbf{N}\) of polynomial functions on \(\operatorname{End}(\mathbf{V})\) .
\item
  If \(x \in \mathfrak{n}\) and \(a \in \mathfrak{a}\) ,
\end{enumerate}

\[
(\exp \operatorname{ad} _ {\mathfrak {a}} x). a = (\exp x) a (\exp (- x)).
\]

\begin{enumerate}
\def\labelenumi{(\roman{enumi})}
\setcounter{enumi}{2}
\tightlist
\item
  Let \(V'\) be a finite dimensional vector space, \(n'\) a Lie subalgebra of \(\mathfrak{gl}(V')\) consisting of nilpotent elements, \(\rho\) a homomorphism from n to \(n'\) . Let \(\pi\) be the map \(\exp x \mapsto \exp \rho(x)\) from \(\exp n\) to \(\exp n'\) . Then \(\pi\) is a group homomorphism.
\end{enumerate}

By Engel's theorem, we can identify V with \(k^n\) in such a way that \(\mathfrak{n}\) is a subalgebra of \(\mathfrak{n}(n,k)\) (the Lie subalgebra of \(\mathbf{M}_n(k)\) consisting of the lower triangular matrices with zeros on the diagonal). For \(s \geq 0\) , let \(\mathfrak{n}_s(n,k)\) be the set of \((x_{ij})_{1 \leq i,j \leq n} \in \mathbf{M}_n(k)\) such that \(x_{ij} = 0\) for \(i - j < s\) . Then

\[
[ \mathfrak {n} _ {s} (n, k), \mathfrak {n} _ {s ^ {\prime}} (n, k) ] \subset \mathfrak {n} _ {s + s ^ {\prime}} (n, k)
\]

(Chap. II, §4, no. 6, Remark), and the Hausdorff series defines a polynomial map \((a,b)\mapsto\mathrm{H}(a,b)\) from \(\mathfrak{n}(n,k)\times\mathfrak{n}(n,k)\) to \(\mathfrak{n}(n,k)\) (Chap. II, §6, no. 5, Remark 3); this map makes \(\mathfrak{n}(n,k)\) into a group (Chap. II, §6, no. 5, Prop. 4). By Chap. II, §6, no. 1, Remark 4, the maps \(x\mapsto\exp x\) from \(\mathfrak{n}(n,k)\) to \(1+\mathfrak{n}(n,k)\) and \(y\mapsto\log y\) from \(1+\mathfrak{n}(n,k)\) to \(\mathfrak{n}(n,k)\) are inverse bijections and are polynomial; by Chap. II, §6, no. 5, Prop. 3, these maps are isomorphisms of groups if \(\mathfrak{n}(n,k)\) is given the group law \((a,b)\mapsto\mathrm{H}(a,b)\) and if \(1+\mathfrak{n}(n,k)\) is considered as a subgroup of \(\mathbf{GL}_{n}(k)\) . Assertions (i) and (iii) of the lemma now follow. Let \(x\in\mathfrak{n}\) . Denote by \(\mathrm{L}_{x},\mathrm{R}_{x}\) the maps \(u\mapsto xu,u\mapsto ux\) from \(\mathfrak{a}\) to \(\mathfrak{a}\) , which commute and are nilpotent. We have \(\mathrm{ad}_{\mathfrak{a}}x=\mathrm{L}_{x}-\mathrm{R}_{x}\) , so, for all \(a\in\mathfrak{a}\) ,

\[
\begin{array}{c} (\exp \operatorname{ad} _ {\mathfrak {a}} x) a = (\exp (\mathrm{L} _ {x} - \mathrm{R} _ {x})) a = (\exp \mathrm{L} _ {x}) (\exp \mathrm{R} _ {- x}) a \\ = \sum_ {i, j \geq 0} \frac {\mathrm{L} _ {x} ^ {i}}{i !} \frac {\mathrm{R} _ {- x} ^ {j}}{j !} a = (\exp x) a (\exp (- x)). \end{array}\tag{1}
\]

With the notation in Lemma 1, \(\pi\) is called the linear representation of \(\exp \mathfrak{n}\) compatible with the given representation \(\rho\) of \(\mathfrak{n}\) on \(V'\) . When \(k\) is \(\mathbf{R}\) , \(\mathbf{C}\) , or a non-discrete complete ultrametric field, \(\rho = L(\pi)\) by the properties of exponential maps (Chap. III, §4, no. 4, Cor. 2 of Prop. 8).

PROPOSITION 1. Let \(\mathfrak{g}\) be a Lie algebra, \(\mathfrak{n}\) a subalgebra of \(\mathfrak{g}\) such that \(\operatorname{ad}_{\mathfrak{g}}x\) is nilpotent for all \(x\in \mathfrak{n}\) . Then \(e^{\operatorname{ad}_{\mathfrak{g}}\mathfrak{n}}\) is a subgroup of \(\operatorname{Aut}_e(\mathfrak{g})\) .

This follows immediately from Lemma 1 (i).

In particular, if \(\mathfrak{n}\) is the nilpotent radical of \(\mathfrak{g}\) , \(e^{\mathrm{ad}_{\mathfrak{g}}\mathfrak{n}}\) is the group of special automorphisms of \(\mathfrak{g}\) (Chap. I, §6, no. 8, Def. 6).

Remarks. 1) Let V be a finite dimensional vector space, g a Lie subalgebra of \(\mathfrak{a} = \mathfrak{gl}(V)\) , x an element of g such that \(ad_{g}x\) is nilpotent. Then there exists a nilpotent element n of a such that \(ad_{a}n\) extends \(ad_{g}x\) . Indeed, let s, n be the semi-simple and nilpotent components of x; then \(ad_{a}s\) and \(ad_{a}n\) are the semi-simple and nilpotent components of \(ad_{a}x\) (Chap. I, §5, no. 4, Lemma 2), so \(ad_{a}s\) and \(ad_{a}n\) leave g stable, and \(ad_{a}s|g\) and \(ad_{a}n|g\) are the semi-simple and nilpotent components of \(ad_{g}x\) ; consequently, \(ad_{g}x = ad_{a}n|g\) , which proves our assertion. In view of Lemma 1 (ii), every elementary automorphism of g extends to an automorphism of a of the form \(u \mapsto mum^{-1}\) where \(m \in \mathbf{GL}(V)\) .

\begin{enumerate}
\def\labelenumi{\arabic{enumi})}
\setcounter{enumi}{1}
\tightlist
\item
  Let V be a finite dimensional vector space. For all \(g \in \mathbf{SL}(V)\) , let \(\varphi(g)\) be the automorphism \(x \mapsto gxg^{-1}\) of \(\mathfrak{gl}(V)\) . Then
\end{enumerate}

\[
\operatorname{Aut} _ {e} (\mathfrak {g l} (V)) = \varphi (\mathbf {S L} (V)).
\]

Indeed, by (1), \(\operatorname{Aut}_{e}(\mathfrak{gl}(V))\) is contained in \(\varphi(\mathbf{SL}(V))\) , and the opposite inclusion follows from Algebra, Chap. III, §8, no. 9, Prop. 17 and (1). An analogous argument shows that \(\operatorname{Aut}_{e}(\mathfrak{sl}(V))\) is the set of restrictions of elements of \(\varphi(\mathbf{SL}(V))\) to \(\mathfrak{sl}(V)\) .

